\documentclass[10pt]{article}
\usepackage[utf8]{inputenc}
\usepackage[T1]{fontenc}
\usepackage{mathtools}
\usepackage{amssymb}
\usepackage{circuitikz}
\usepackage{fontspec}
\usepackage{array}
\usepackage{float}

\setmainfont{Times New Roman}

%
% NOTE: schematics generated using the online tool: 'https://www.circuit2tikz.tf.fau.de/designer'.
% Special thanks to them!
%

\title{Analysis of the ARP Quadra Phaser}
\author{QuBi}
\date{\today}

\begin{document}

\maketitle


\section*{Foreword}

A few methodological points for approaching such a problem correctly:

\begin{itemize}
  \item \textbf{A synthesizer module is not a laboratory instrument.}
  It is not designed to be perfect or to precisely implement the function it advertises.
  It may simply be an approximation that \textit{sounds good} --- and ultimately, that is the only thing a synthesizer is asked to do.
  As such, when we speak of a ``filter'' or a ``phaser'', there are implicit fine-print conditions that are not always spelled out. For example:
  \begin{itemize}
    \item Ladder Filter: only operates correctly at very small amplitudes; noticeable gain loss at high Q, etc.
    \item VCO: waveform is a rough approximation, especially at extreme frequencies.
  \end{itemize}
  The limitations and deviations from ideal behaviour are never stated explicitly.
  This can introduce assumptions in the analysis that turn out to be misleading.

  \item \textbf{A technical service schematic is not designed to be pedagogical.}
  Such documents are primarily intended for assemblers, repair technicians, and the like.
  The draughtsmen therefore had different priorities:
  \begin{itemize}
    \item fitting a coherent set of components onto a single page,
    \item making it easy to locate individual components.
  \end{itemize}
  At the time, schematics were drawn by hand, leaving little room for revision.
  As a result, what one sees on the page is not necessarily the most conducive representation for understanding.
  Do not hesitate to redraw or rearrange the schematic to better grasp its origin and operating principle.

  \item \textbf{A purely computational analytical approach rarely helps in understanding how a system works.}
  It is powerful for quantitative analysis (how much? what is the exact relationship between two quantities? etc.),
  but it tends to obscure the overarching principles that explain why the system performs a given function.
  Setting up equations should only be done for a specific sub-system, and only to shed light on a very precise aspect.

  \item \textbf{There is no point overcomplicating matters by considering fifteenth-order effects.}
  In most cases, these circuits were designed using first-order models.
  There is therefore no value in burdening the analysis with unnecessary assumptions.
\end{itemize}


\section*{Analysis}

\subsection*{Reducing the Problem}

The repetitive ladder structure immediately suggests focusing, as a first step, on the role of a single rung.
By induction, this should provide good insight into the overall behaviour.

The functions of the surrounding stages are abstracted away, as they are secondary:
\begin{itemize}
  \item the lower section sinks a current $I$ in each branch,
  \item the upper section injects the signal to be phase-shifted.
\end{itemize}

Since we are working in the frequency domain (harmonic analysis), capacitors can be replaced by simple resistors to fix ideas.

The circuit exhibits a manifest symmetry: the left half is a mirror image of the right half.
Again, this invites us to consider only one half of the circuit.

This kind of symmetry is quite common: it is routinely used to separate the operating point
(static and identical in each branch) from the signal component (varying and of opposite sign in each branch).

A simple subtraction then eliminates the unwanted operating point.


\subsection*{Qualitative Operation}

We observe the voltage at the transistor's emitter.

Let us consider the extreme values of the control current.
By continuity, we will infer the general behaviour.

\paragraph{Zero control current.}
If the current is zero, the transistor is inactive and only the resistor connects input to output.
With no current, the collector voltage is reproduced identically at the emitter: the stage is transparent.

\paragraph{Increasing the control current.}
The collector potential is imposed by a voltage source and therefore remains essentially unchanged.
This current flows through $R_C$, creating a voltage drop between collector and emitter that tends to pull the emitter potential down
(which is possible since the current sink below does not, by definition, impose a voltage).
As the base voltage is fixed, $v_{BE}$ increases and begins driving a current through $Q$.

$Q$ and $R_C$ share the same constant total current $I_\mathrm{cmd}$.
Consequently, the onset of conduction in $Q$ alters the distribution of current between $Q$ and $R_C$,
and therefore the value of $v_{BE}$.

A negative feedback loop on $I_C$ takes shape:
\[
  I_{R_C} \uparrow \;\Rightarrow\; v_{BE} \uparrow \;\Rightarrow\; I_C \uparrow \;\Rightarrow\; I_{R_C} \downarrow
\]

\begin{quote}
  \textbf{Stability aside.}
  This type of mechanism typically converges naturally toward equilibrium, but it can also oscillate.
  A precise knowledge of the loop parameters (gain and phase shift) would be needed to draw a firm conclusion; however, in the present case:
  \begin{itemize}
    \item there are no significant delay elements (e.g.\ capacitors) that would give cause for doubt,
    \item simulations and bench tests support the stability hypothesis.
  \end{itemize}
  Delay is an important stability criterion.
  It is somewhat analogous to the situation where two people walking toward each other on a pavement both step aside --- but sometimes their reaction times are such that they mirror each other, and the awkward dance continues for a while.
  Here, stability can reasonably be assumed.
\end{quote}

An equilibrium is thus established between $I_C$ and $I_{R_C}$, and it depends on the control current:
a higher $I_\mathrm{cmd}$ leads to a higher $v_{BE}$.

\paragraph{Effect on the input signals.}
From the emitter's perspective, the upper input is still seen through resistor $R_C$.
The second input, however, is seen through a dynamic resistance that decreases as the control current increases.
(In this region of the curve, a very small $\Delta V$ produces a very large $\Delta I$, so $r_d = \Delta V / \Delta I$ is very small.)

Since the $v_{BE}$ vs.\ $I_C$ characteristic is exponential, a point will be reached where \ldots


\subsection*{Application to the Complete Schematic}

As noted above, the differential (symmetrised) configuration eliminates the bias operating point.

The operation can then be summarised in the following key points:
\begin{enumerate}
  \item Each stage performs a blend between the signal and its $180°$-phase-shifted version.
  \item The mixing ratio depends on the control current.
  \item Replacing the resistors with capacitors makes the ratio frequency-dependent:
    \begin{itemize}
      \item $f \ll f_0$: the original signal passes through,
      \item $f \gg f_0$: the $180°$-shifted version passes through,
      \item $f \approx f_0$: \ldots
    \end{itemize}
\end{enumerate}

\end{document}